% Experimental methodology section for OrbitGraph paper.

\section{Experimental Methodology}
\label{sec:experiments}

We evaluate OrbitGraph against a distributed OSPF baseline under identical
emulated constellations, traffic, and failure scripts.
All runs use StarryNet's full experiment profile on a single host; results are
aggregated across repeated, seeded trials and reported with mean and standard
deviation where applicable.

\subsection{Constellation scenarios and scaling grid}

We study six square grids from $5{\times}5$ through $10{\times}10$ orbital
planes and satellites per plane.
Each grid comprises $O \cdot S$ satellites plus two ground stations ($N = O
\cdot S + 2$ nodes total), ranging from 27 to 102 nodes.
Grids share the Starlink-inspired parameters in Section~\ref{sec:simulation}
(550\,km altitude, 53$^\circ$ inclination, plus-grid ISLs, LeastDelay policy,
instant GSL handover).
Ground-station endpoints for measurement are derived automatically as nodes
$O \cdot S + 1$ and $O \cdot S + 2$.

Emulation duration is sized per grid so that at least one GSL handover occurs
and a late steady-state sample is captured before teardown:
100\,s ($5{\times}5$), 120\,s ($6{\times}6$), 250\,s ($7{\times}7$),
400\,s ($8{\times}8$ and $10{\times}10$), and 350\,s ($9{\times}9$).
These values follow our paper-scale plan (\texttt{simulation\_paper.json}) and
ensure handover-relative probes fall inside the emulation window for every size.

\subsection{Compared routing modes}

We run two modes on each grid and seed:

\begin{itemize}
  \item \textbf{OSPF baseline.} The BIRD routing daemon runs OSPF in every
        container; the emulation waits for OSPF convergence after initialization
        and reconverges after each topology mutation.
  \item \textbf{OrbitGraph (SDN).} A centralized controller recomputes
        delay-weighted shortest paths (Dijkstra) on each routing event,
        pushes only changed kernel routes (\texttt{ip route replace} before
        \texttt{ip route del} for make-before-break), and performs
        \emph{proactive handover}: post-handover FIB entries are installed
        after new GSLs are added but before old GSLs are removed.
\end{itemize}

All other emulation steps---delay updates every 10\,s, scripted link damage and
recovery, and ping scheduling---are identical between modes.
Only the control-plane path differs, yielding a paired comparison on the same
dynamic graph sequence.

\subsection{Workload and event script}

Each run follows the \emph{full} profile:

\begin{enumerate}
  \item \textbf{Initialization.} Containers and links are created from second-1
        delay matrices; routing is installed (OSPF convergence or OrbitGraph
        initial FIB push). A ground-station-to-ground-station ping
        (\texttt{post\_init} phase) verifies baseline reachability.
  \item \textbf{Link damage ($t{=}5$).} Thirty percent of satellites are chosen
        pseudorandomly (seeded per repetition); all interfaces on those nodes
        receive 100\% loss via \texttt{netem}, modeling satellite failure.
  \item \textbf{Link recovery ($t{=}10$).} The same interfaces revert to the
        configured 1\% loss rate.
  \item \textbf{GSL handover.} The first entry in \texttt{Topo\_leo\_change.txt}
        adds a new ground--satellite link and removes an old one at a
        geometry-dependent tick (e.g., $t{=}53$ on $5{\times}5$, $t{=}23$ on
        $6{\times}6$). Pings are scheduled at the handover instant, two seconds
        later (\texttt{post\_handover}), and near emulation end
        (\texttt{steady}), so phases align across grid sizes despite different
        handover times.
\end{enumerate}

\subsection{Metrics}

\paragraph{Data plane.}
The primary outcome is \textbf{data-plane outage duration} (\texttt{outage\_ms})
around each routing event.
A continuous \texttt{ping -D -O} probe (10\,Hz) runs inside the source ground-
station container for the entire emulation; timestamps are aligned to control-
plane snapshot wall times.
Outage onset requires a sustained loss burst ($\geq$3 consecutive unreplied
ICMP sequences); recovery requires $\geq$3 consecutive replies.
This metric is symmetric for OSPF and OrbitGraph and captures multi-second OSPF
black holes that tick-scheduled pings would miss when route install blocks the
emulation loop.

We additionally report \textbf{ping loss (\%)} and \textbf{mean RTT (ms)} per
phase (\texttt{post\_init}, \texttt{handover}, \texttt{post\_handover},
\texttt{steady}) for the ground-station pair, and qualitative path checks via
traceroute artifacts where applicable.

\paragraph{Control plane.}
For OrbitGraph we record \texttt{compute\_ms} (graph load plus Dijkstra),
\texttt{install\_ms} (kernel route push), \texttt{nodes\_touched}, and route
counts (\texttt{installed}, \texttt{deleted}, \texttt{failed}) per snapshot
reason (\texttt{init}, \texttt{damage\_recovery}, \texttt{topology\_change},
\texttt{proactive\_handover}).
For OSPF we record BIRD route-dump \texttt{collection\_ms} as an observation
cost; it is not equated with OrbitGraph install time.
Headline control-plane comparisons at handover contrast OSPF collection against
OrbitGraph compute plus incremental install at the \texttt{proactive\_handover}
snapshot.

\subsection{Repetitions and aggregation}

We execute \textbf{ten independent repetitions} per (grid, mode) pair
(120 runs total across the scaling study).
Repetition $i$ uses RNG seed $1000 + i$, fixing the damaged-satellite subset
while preserving natural Docker and kernel timing variation.
Runs are orchestrated by \texttt{experiments/compare\_batch.py} with
\texttt{--profile full} and merged into summary CSVs
(\texttt{outage\_summary.csv}, \texttt{ping\_summary\_by\_phase.csv},
\texttt{control\_summary\_by\_reason.csv}) for plotting.
Each batch checks for expected phases and routing-event reasons and flags
incomplete samples.

Figures in this paper aggregate the first-handover outage and handover-phase
loss/RTT unless noted otherwise; cross-size trends use the \texttt{nodes}
column ($N = O \cdot S + 2$).

\subsection{Execution environment}

Experiments ran on a single Apple M2 Pro workstation (10 CPU cores, 32\,GiB
RAM) under macOS~14.6 with OrbStack providing the Linux container runtime
(Docker API 29.x, 19.6\,GiB allocated to the VM).
StarryNet \texttt{local\_mode} executes \texttt{docker exec} and \texttt{tc}
commands on the host without remote SSH.
Python~3.11 drives batch orchestration; each repetition tears down
\texttt{ovs\_container\_*} instances and associated bridge networks before the
next seed starts.

This setup reflects a realistic researcher workstation deployment of
StarryNet~\cite{lai2023starrynet} and keeps OSPF and OrbitGraph runs colocated
on the same hardware, container image, and emulation timeline.
